\documentclass[11pt,a4paper,spanish]{book}
\usepackage{estilo_unir-1}
\usepackage{apacite}
\usepackage{hyperref}
\usepackage{longtable}
\usepackage{booktabs}
\usepackage{array}
\usepackage{tabularx}
\numberwithin{equation}{chapter}
\numberwithin{figure}{chapter}

%---------------------------
% METADATOS
%---------------------------
\title{ProcureWatch Analytics: Plataforma Multiagente para la Detecci{\'o}n de Patrones
de Riesgo en Contrataci{\'o}n P{\'u}blica mediante Open Data, An{\'a}lisis de Redes,
Procesamiento de Lenguaje Natural y Modelos Locales de Inteligencia Artificial}
\titulacion{M{\'a}ster Universitario en An{\'a}lisis y Visualizaci{\'o}n de Datos Masivos}
\author{Sebasti{\'a}n Alfonso Correa y Satur{\'i}a Mar{\'i}a Hurtado {\'A}lvarez}
\date{22 de abril de 2026}
\director{Marlon C{\'a}rdenas Bonett}
\nombreciudad{Madrid}

\begin{document}
\renewcommand{\listfigurename}{\'Indice de Ilustraciones}
\renewcommand{\listtablename}{\'Indice de Tablas}
\renewcommand{\contentsname}{\'Indice de Contenidos}
\renewcommand{\figurename}{Figura}
\renewcommand{\tablename}{Tabla}

\maketitle
\frontmatter
\tableofcontents
\listoffigures
\listoftables

%=========================================================
\chapter{Resumen}
%=========================================================

La contrataci{\'o}n p{\'u}blica en Espa{\~n}a moviliza anualmente entre el 12 y el 15\,\% del
Producto Interior Bruto, sin contar los contratos menores que no exigen publicidad,
y su gesti{\'o}n est{\'a} expuesta a riesgos documentados de fraccionamiento, colusi{\'o}n y
adjudicaci{\'o}n dirigida. A pesar de la disponibilidad de datos abiertos en formato
OCDS desde 2007, no existe actualmente ninguna plataforma open-source, reproducible
y validada que combine, sobre datos espa{\~n}oles, indicadores de red flags, an{\'a}lisis
de redes societarias y procesamiento de lenguaje natural sobre documentos contractuales.

Este trabajo propone \textbf{ProcureWatch Analytics}, una plataforma anal{\'i}tica
multiagente que cubre ese vac{\'i}o. El sistema integra: un \textit{pipeline} de
adquisici{\'o}n e integraci{\'o}n de datos OCDS; un motor de m{\'a}s de 15 indicadores de
red flags con modelos de machine learning (Isolation Forest, Positive-Unlabeled
Learning y Random Forest); un m{\'o}dulo de an{\'a}lisis de redes basado en el algoritmo
Leiden para la detecci{\'o}n de comunidades en el grafo bipartito empresa-organismo-contrato;
un m{\'o}dulo NLP con OCR (Docling), reconocimiento de entidades (spaCy) y recuperaci{\'o}n
aumentada con generaci{\'o}n (RAG h{\'i}brido sobre Qdrant); y un \textit{front-end} anal{\'i}tico
con scoring de riesgo explicable mediante SHAP y LIME.

El prototipo se valida sobre la divisi{\'o}n CPV~71 (Servicios de arquitectura,
ingenier{\'i}a e inspecci{\'o}n), seleccionada por sus propiedades metodol{\'o}gicas
favorables para la evaluaci{\'o}n de los tres m{\'o}dulos, procesando el \textit{dataset}
longitudinal del BOE 2014--2024 \cite{MunozPla2026} (aprox.\ 97.000 contratos) y datos
de la Plataforma de Contrataci{\'o}n del Sector P{\'u}blico (PLACE). Los componentes se
eval{\'u}an con m{\'e}tricas est{\'a}ndar: precisi{\'o}n de red flags ($>70$\,\%), modularidad
de grafos ($Q>0{,}30$), calidad RAG (Faithfulness $>0{,}80$) y usabilidad del panel
(SUS\,$>68$), generando fichas de riesgo explicables que apoyan la priorizaci{\'o}n de
casos para revisi{\'o}n por organismos de control.

\medskip
{\bf Palabras clave:} contrataci{\'o}n p{\'u}blica, red flags, detecci{\'o}n de irregularidades,
machine learning, an{\'a}lisis de redes, NLP, datos abiertos, CPV~71, OCDS, scoring de
riesgo, plataforma multiagente, explicabilidad.

%=========================================================
\chapter{Abstract}
%=========================================================

Public procurement in Spain represents between 12 and 15\,\% of Gross Domestic Product
annually and is systematically exposed to documented risks of contract splitting,
collusion, and directed award. Despite the availability of structured open data in OCDS
format since 2007, no open-source, reproducible and validated platform currently exists
that combines, on Spanish data, red flag indicators, corporate network analysis and
natural language processing over contracting documents.

This work proposes \textbf{ProcureWatch Analytics}, a multi-agent analytical platform
that fills this gap. The system integrates: an OCDS data acquisition and integration
pipeline; an engine of 15+ red flag indicators with machine learning models (Isolation
Forest, Positive-Unlabeled Learning and Random Forest); a network analysis module based
on the Leiden algorithm for community detection in the bipartite company-organisation-contract
graph; an NLP module with OCR (Docling), named entity recognition (spaCy) and hybrid
retrieval-augmented generation (RAG over Qdrant); and an analytical front-end with
explainable risk scoring via SHAP and LIME.

The prototype is validated on CPV division 71 (Architecture, Engineering and Inspection
Services), selected for its favourable methodological properties for evaluating all three
modules, processing the longitudinal BOE 2014--2024 dataset \cite{MunozPla2026}
(approx.\ 97,000 contracts) and PLACE data. Components are evaluated using standard
metrics: red flag precision ($>70$\,\%), graph modularity ($Q>0.30$), RAG quality
(Faithfulness $>0.80$) and dashboard usability (SUS\,$>68$), generating explainable
risk reports to support case prioritisation for review by oversight bodies.

\medskip
{\bf Keywords:} public procurement, red flags, irregularity detection, machine learning,
network analysis, NLP, open data, CPV~71, OCDS, risk scoring, multi-agent platform,
explainability.

%=========================================================
\mainmatter
\chapter{Introducci{\'o}n}
%=========================================================

\section{Motivaci{\'o}n y justificaci{\'o}n}

La contrataci{\'o}n p{\'u}blica es uno de los mecanismos de gasto p{\'u}blico con mayor
exposici{\'o}n al riesgo de irregularidad. En Espa{\~n}a representa entre el 12 y el
15\,\% del PIB nacional \cite{OECD2016Corruption} y moviliza anualmente decenas de
miles de millones de euros en obras, servicios y suministros gestionados por m{\'a}s de
10.000 {\'o}rganos de contrataci{\'o}n activos. Esta magnitud, combinada con la
fragmentaci{\'o}n administrativa del sistema, la asimetr{\'i}a de informaci{\'o}n entre
licitadores y organismos, y la complejidad t{\'e}cnica de los procedimientos, convierte
la contrataci{\'o}n p{\'u}blica en un {\'a}mbito especialmente propicio para la aparici{\'o}n
de comportamientos irregulares \cite{OECD2016Corruption,OCP2024RedFlags}.

Desde 2007, la legislaci{\'o}n espa{\~n}ola obliga a publicar los contratos en formatos
abiertos y reutilizables \cite{LCSP2017,Ley192013}, lo que ha generado una infraestructura
de datos sin precedentes. El reciente \textit{dataset} longitudinal del BOE 2014--2024
\cite{MunozPla2026}, con aproximadamente 97.000 contratos estructurados seg{\'u}n el
est{\'a}ndar Open Contracting Data Standard (OCDS), y los m{\'a}s de 200.000 procedimientos
anuales publicados en la Plataforma de Contrataci{\'o}n del Sector P{\'u}blico (PLACE)
\cite{OIReScon2025} representan una oportunidad extraordinaria para el an{\'a}lisis
sistem{\'a}tico mediante t{\'e}cnicas avanzadas de datos masivos.

La literatura cient{\'i}fica internacional ha demostrado de forma consistente que los
m{\'e}todos de machine learning y an{\'a}lisis de redes mejoran significativamente la
detecci{\'o}n de irregularidades respecto a los enfoques basados {\'u}nicamente en reglas
expl{\'i}citas o auditor{\'i}a manual \cite{Wachs2025ML,Fazekas2020DIGIWHIST,Coviello2023Italian}.
Sin embargo, este conocimiento acumulado no se ha traducido en herramientas integradas,
reproducibles y aplicadas al caso espa{\~n}ol. Los estudios de mayor impacto se han
realizado sobre datos italianos \cite{Coviello2023Italian}, portugueses
\cite{Martins2022Network}, mexicanos \cite{Wachs2025ML} o centroeuropeos
\cite{Fazekas2020DIGIWHIST}, sin que exista un prototipo equivalente para Espa{\~n}a.

Esta doble brecha ---disponibilidad de datos de alta calidad sin explotaci{\'o}n
sistem{\'a}tica, y conocimiento metodol{\'o}gico sin implementaci{\'o}n local--- constituye la
justificaci{\'o}n central de este trabajo.

\section{Preguntas de investigaci{\'o}n}

A partir del an{\'a}lisis del contexto desarrollado en el Cap{\'i}tulo~\ref{chap:contexto}, el
presente trabajo aborda cuatro preguntas de investigaci{\'o}n:

\begin{description}
  \item[PI-1] \textit{{\textquestiondown}Cu{\'a}l es el volumen real de contratos publicados en
    PLACE y cu{\'a}ntos organismos participan, y qu{\'e} implicaciones tiene esta escala para el
    dise{\~n}o de sistemas de supervisi{\'o}n?}

  \item[PI-2] \textit{{\textquestiondown}Qu{\'e} tipos de irregularidades en contrataci{\'o}n p{\'u}blica
    se han documentado en Espa{\~n}a en los {\'u}ltimos diez a{\~n}os, y qu{\'e} patrones de datos
    permiten detectarlas?}

  \item[PI-3] \textit{{\textquestiondown}Cu{\'a}les son las limitaciones estructurales de los sistemas
    actuales de auditor{\'i}a y control, y qu{\'e} tipos de an{\'a}lisis quedan fuera de su alcance?}

  \item[PI-4] \textit{{\textquestiondown}Qu{\'e} evidencia emp{\'i}rica existe sobre la eficacia del
    an{\'a}lisis de datos masivos, el machine learning y el an{\'a}lisis de redes para detectar
    patrones de riesgo en contrataci{\'o}n p{\'u}blica?}
\end{description}

\section{Planteamiento del trabajo}

Este TFM propone dise{\~n}ar, desarrollar y evaluar \textbf{ProcureWatch Analytics}, una
plataforma multiagente que, a partir de datos abiertos de contrataci{\'o}n p{\'u}blica
espa{\~n}ola, detecta y visualiza patrones de riesgo mediante cinco componentes integrados:

\begin{enumerate}
  \item Un \textit{pipeline} de adquisici{\'o}n, limpieza, normalizaci{\'o}n y enriquecimiento
    de datos Open Data en esquema OCDS unificado.
  \item Un motor anal{\'i}tico de red flags con m{\'a}s de 15 indicadores basados en literatura
    y tres enfoques de machine learning: reglas expl{\'i}citas, Isolation Forest y
    Positive-Unlabeled Learning.
  \item Un m{\'o}dulo de an{\'a}lisis de redes societarias sobre el grafo bipartito
    empresa-organismo-contrato, con detecci{\'o}n de comunidades mediante el algoritmo Leiden.
  \item Un sistema de procesamiento de documentos contractuales (pliegos, resoluciones,
    memorias justificativas) mediante OCR, reconocimiento de entidades y motor RAG h{\'i}brido.
  \item Un \textit{front-end} anal{\'i}tico con grafo interactivo, panel de indicadores de
    riesgo, evoluci{\'o}n temporal y fichas de riesgo explicables.
\end{enumerate}

El prototipo se valida sobre el dominio CPV~71 (Servicios de arquitectura, ingenier{\'i}a
e inspecci{\'o}n), cuya elecci{\'o}n se fundamenta metodol{\'o}gicamente en la
Secci{\'o}n~\ref{sec:dominio_cpv71}. La metodolog{\'i}a de desarrollo detallada, incluyendo
fases, cronograma de 18~semanas, organizaci{\'o}n del equipo y entregables por fase, se
documenta en el documento paralelo \texttt{metodologia\_tfm.pdf}.

\section{Estructura del trabajo}

\begin{description}
  \item[Cap{\'i}tulo~2 --- An{\'a}lisis del contexto.] Responde a las cuatro preguntas de
    investigaci{\'o}n con datos justificativos, caracteriza el dominio CPV~71 seleccionado y
    posiciona el proyecto frente a la brecha identificada en la literatura.
  \item[Cap{\'i}tulo~3 --- Estado del arte.] Revisi{\'o}n cr{\'i}tica de la literatura en cinco
    bloques tem{\'a}ticos: datos abiertos y OCDS; red flags, indicadores de riesgo y
    machine learning; an{\'a}lisis de redes; NLP y procesamiento documental; visualizaci{\'o}n
    anal{\'i}tica y explicabilidad.
  \item[Cap{\'i}tulo~4 --- Objetivos.] Objetivo general SMART y ocho objetivos espec{\'i}ficos
    medibles y operativos, cada uno con m{\'e}tricas de {\'e}xito verificables.
  \item[Cap{\'i}tulo~5 --- Marco normativo.] Normativa relevante en el dise{\~n}o del sistema.
  \item[Cap{\'i}tulos~6--8 ---] Desarrollo t{\'e}cnico, resultados, evaluaci{\'o}n y conclusiones
    (a desarrollar en la fase de implementaci{\'o}n).
\end{description}

%=========================================================
\chapter{An{\'a}lisis del Contexto}
\label{chap:contexto}
%=========================================================

\section{Dimensi{\'o}n econ{\'o}mica e institucional de la contrataci{\'o}n p{\'u}blica en Espa{\~n}a}

La contrataci{\'o}n p{\'u}blica en Espa{\~n}a constituye uno de los instrumentos de pol{\'i}tica
econ{\'o}mica de mayor impacto directo: representa entre el 12 y el 15\,\% del PIB
nacional \cite{OECD2016Corruption} y canaliza la inversi{\'o}n p{\'u}blica en infraestructura,
servicios esenciales, tecnolog{\'i}a y asistencia t{\'e}cnica. Su gesti{\'o}n est{\'a} distribuida
entre m{\'a}s de 10.000 {\'o}rganos de contrataci{\'o}n activos que abarcan todos los niveles
administrativos: la Administraci{\'o}n General del Estado, 17 comunidades aut{\'o}nomas, 50
diputaciones provinciales, m{\'a}s de 8.000 municipios, universidades p{\'u}blicas y el
amplio sector p{\'u}blico institucional \cite{MinHacienda2024PLACE}.

Esta fragmentaci{\'o}n introduce una tensi{\'o}n estructural: por un lado, los principios de
concurrencia, publicidad y transparencia exigen que la informaci{\'o}n sea p{\'u}blica y
accesible; por otro, la dispersi{\'o}n entre cientos de plataformas y formatos heterog{\'e}neos
dificulta el an{\'a}lisis sistem{\'a}tico. La Tabla~\ref{tab:volumen_datos} resume el volumen
y la cobertura de las principales fuentes de datos disponibles.

\begin{table}[ht]
\centering
\small
\caption{Fuentes y volumen de datos disponibles para contrataci{\'o}n p{\'u}blica en Espa{\~n}a y Europa}
\label{tab:volumen_datos}
\begin{tabularx}{\textwidth}{Xrr}
\toprule
\textbf{Fuente / indicador} & \textbf{Volumen} & \textbf{Referencia} \\
\midrule
Licitaciones anuales PLACE (2024, sin contratos menores) & 206.242 & \citeA{OIReScon2025} \\
{\'O}rganos de contrataci{\'o}n activos en PLACE & $>$10.000 & \citeA{MinHacienda2024PLACE} \\
Contratos BOE 2014--2024 en formato OCDS & $\sim$97.000 & \citeA{MunozPla2026} \\
Contratos CPV~71 en \textit{dataset} BOE & 3.443 & \citeA{MunozPla2026} \\
Importe total adjudicado CPV~71 & $>$2.179\,M\texteuro & \citeA{MunozPla2026} \\
Empresas adjudicatarias distintas CPV~71 & 2.031 & \citeA{MunozPla2026} \\
Contratos europeos en OpenTender & $>$19 millones & \citeA{OpenTenderEU} \\
Indicadores de red flags (OCP + DIGIWHIST) & $>$150 & \citeA{OCP2024RedFlags} \\
\bottomrule
\end{tabularx}
\end{table}

\section{Escala del sistema de contrataci{\'o}n y sus implicaciones analíticas (PI-1)}
\label{sec:pi1}

La Plataforma de Contrataci{\'o}n del Sector P{\'u}blico (PLACSP) centraliza la publicaci{\'o}n
de licitaciones y adjudicaciones de toda la administraci{\'o}n espa{\~n}ola, y constituye la
fuente primaria de datos abiertos del sistema \cite{MinHacienda2024PLACE}. Seg{\'u}n los
informes anuales de la Oficina Independiente de Regulaci{\'o}n y Supervisi{\'o}n de la
Contrataci{\'o}n (OIReScon), en el ejercicio 2024 se registraron \textbf{206.242 licitaciones},
sin incluir contratos menores, lo que representa una parte sustancial del total de la
actividad contractual gestionada en Espa{\~n}a \cite{OIReScon2025}. Si se incorporan
contratos menores, adjudicaciones directas y desgloses por lotes, el volumen total de
registros anuales asciende a varios millones, aunque la cifra var{\'i}a seg{\'u}n el criterio
de contabilizaci{\'o}n (contrato, expediente o lote) \cite{OIReScon2025}.

Esta escala tiene una implicaci{\'o}n anal{\'i}tica directa: ning{\'u}n equipo de auditores
puede revisar manualmente el conjunto de contratos. La supervisi{\'o}n exhaustiva requiere
necesariamente herramientas autom{\'a}ticas de priorizaci{\'o}n de riesgos. Adem{\'a}s, la
fragmentaci{\'o}n de los datos entre m{\'u}ltiples administraciones introduce heterogeneidad
en formatos, campos disponibles y calidad de la informaci{\'o}n publicada \cite{OCP2023EUOpenData},
lo que exige un \textit{pipeline} de normalizaci{\'o}n robusto como condici{\'o}n previa a
cualquier an{\'a}lisis comparativo.

\section{Irregularidades documentadas en contrataci{\'o}n p{\'u}blica espa{\~n}ola (PI-2)}
\label{sec:pi2}

La literatura cient{\'i}fica y los informes de supervisi{\'o}n han identificado y documentado
un conjunto estructurado de tipolog{\'i}as de irregularidad en la contrataci{\'o}n p{\'u}blica
espa{\~n}ola durante la {\'u}ltima d{\'e}cada. A efectos del presente trabajo, se distinguen
dos niveles: irregularidades detectables mediante indicadores tabulares sobre datos OCDS,
e irregularidades que requieren an{\'a}lisis relacional o documental.

\subsection{Irregularidades detectables con indicadores tabulares}

\begin{description}
  \item[Fraccionamiento indebido de contratos.] Divisi{\'o}n artificial del objeto
    contractual para eludir los umbrales de publicidad y concurrencia establecidos por la
    LCSP. Es el indicador con mayor soporte emp{\'i}rico en la literatura sobre contrataci{\'o}n
    espa{\~n}ola \cite{OCP2024RedFlags,OCP2016RedFlags}. En el dominio CPV~71, el
    fraccionamiento es especialmente frecuente en servicios de asistencia t{\'e}cnica
    continuada, donde la partici{\'o}n por fases temporales puede estar o no justificada.

  \item[Uso inadecuado del procedimiento negociado sin publicidad.] Recurso sistem{\'a}tico
    al contrato menor o al procedimiento negociado ---que no exige concurrencia de
    ofertas--- con el mismo proveedor en importes recurrentemente pr{\'o}ximos al umbral
    legal. Los informes OIReScon han documentado un incremento significativo de este tipo
    de irregularidades en los {\'u}ltimos ejercicios \cite{OIReScon2025}.

  \item[Desviaci{\'o}n sistem{\'a}tica entre importe estimado y adjudicado.] Diferencias
    persistentemente elevadas entre el valor estimado de licitaci{\'o}n y el importe real
    adjudicado, lo que puede indicar estimaciones deliberadamente infladas para reducir
    la concurrencia, o presupuestos artificialmente bajos para favorecer la adjudicaci{\'o}n
    a un proveedor concreto \cite{Coviello2023Italian}. En el dataset BOE, diferencias
    superiores al 20\,\% aparecen en el 5--8\,\% de los contratos.

  \item[Plazos de resoluci{\'o}n inusualmente cortos.] Publicaci{\'o}n y adjudicaci{\'o}n en
    per{\'i}odos muy reducidos, o en fechas de baja visibilidad (v{\'i}speras de festivos,
    cierres de ejercicio presupuestario), que pueden limitar la concurrencia efectiva
    \cite{OCP2024RedFlags}.
\end{description}

\subsection{Irregularidades detectables mediante an{\'a}lisis relacional}

\begin{description}
  \item[Concentraci{\'o}n at{\'i}pica de adjudicaciones.] En algunos organismos, un porcentaje
    superior al 60\,\% de los contratos adjudicados en un per{\'i}odo se concentra en dos o
    tres proveedores, patr{\'o}n que excede la variaci{\'o}n esperable en mercados competitivos.
    \citeA{FalconCortes2021} demuestran que la concentraci{\'o}n de adjudicaciones y la baja
    concurrencia son los indicadores con mayor poder predictivo de comportamiento irregular.

  \item[Pr{\'a}cticas colusorias.] Empresas que licitan formalmente como competidores pero
    comparten administradores, domicilios sociales, estructura accionarial o participan
    conjuntamente en UTEs en otros procedimientos, lo que puede indicar coordinaci{\'o}n en
    las ofertas \cite{OCP2024RedFlags}. Estas relaciones s{\'o}lo son detectables mediante
    an{\'a}lisis de redes societarias.

  \item[Adjudicaciones dirigidas con conflicto de inter{\'e}s.] Un caso emblem{\'a}tico en
    Espa{\~n}a es el detectado en Cantabria en 2023 por la Agencia Tributaria: se
    identificaron adjudicaciones sistem{\'a}ticamente dirigidas a empresas vinculadas a un
    funcionario responsable de los informes t{\'e}cnicos de valoraci{\'o}n de las ofertas
    \cite{AEATCantabria2023}. La AEAT se{\~n}al{\'o} expl{\'i}citamente el cruce entre datos
    tributarios y contractuales como l{\'i}nea prioritaria de detecci{\'o}n.
\end{description}

\subsection{Irregularidades detectables mediante an{\'a}lisis documental (NLP)}

Los pliegos de prescripciones t{\'e}cnicas, las memorias justificativas y las resoluciones
de adjudicaci{\'o}n contienen se{\~n}ales textuales que no son capturables por indicadores
tabulares: criterios de adjudicaci{\'o}n redactados con especificidad excepcional para
favorecer a un proveedor concreto (\textit{tailored tenders}), requisitos de experiencia
en proyectos previos del mismo organismo que excluyen a la mayor parte del mercado, o
memorias justificativas de urgencia que no se corresponden con la naturaleza del servicio.
En el dominio CPV~71, la literatura estima que estos criterios subjetivos pueden
representar hasta el 49\,\% de la puntuaci{\'o}n total en procedimientos abiertos con m{\'u}ltiples
criterios \cite{NLPLegalSurvey2024}.

\section{Limitaciones estructurales de los sistemas de control actuales (PI-3)}
\label{sec:pi3}

Los sistemas de auditor{\'i}a y control de la contrataci{\'o}n p{\'u}blica presentan limitaciones
estructurales que justifican el desarrollo de herramientas anal{\'i}ticas avanzadas. Estas
limitaciones se organizan en tres dimensiones:

\subsection{Limitaciones de escala y recursos}

El volumen de contratos publicados anualmente ---m{\'a}s de 200.000 licitaciones, a las que
deben sumarse los contratos menores no sujetos a publicidad obligatoria--- excede
ampliamente la capacidad de revisi{\'o}n manual de los {\'o}rganos de control. La OIReScon,
principal organismo de supervisi{\'o}n, dispone de recursos significativamente inferiores
al volumen de actividad que debe supervisar \cite{OIReScon2025}. Esta desproporci{\'o}n
hace que, en la pr{\'a}ctica, la revisi{\'o}n se realice sobre muestras no representativas o
a partir de denuncias.

\subsection{Limitaciones de datos e integraci{\'o}n}

Los datos de contrataci{\'o}n est{\'a}n distribuidos entre m{\'u}ltiples plataformas y niveles
administrativos con formatos heterog{\'e}neos. El informe de la Open Contracting Partnership
sobre apertura de datos en la UE \cite{OCP2023EUOpenData} constata que, en Espa{\~n}a,
campos cr{\'i}ticos como el NIF del adjudicatario, el importe exacto adjudicado o las fechas
de los hitos contractuales presentan tasas de completitud irregulares. Esta situaci{\'o}n
impide el an{\'a}lisis sistem{\'a}tico y comparativo a escala nacional.

\subsection{Limitaciones metodol{\'o}gicas: car{\'a}cter reactivo y ausencia de perspectiva
relacional y documental}

El enfoque predominante en los sistemas de control es reactivo: act{\'u}an una vez detectada
o denunciada la irregularidad, no antes \cite{OECD2016Corruption}. Adem{\'a}s, los controles
tradicionales presentan dos puntos ciegos espec{\'i}ficos:

\begin{itemize}
  \item \textbf{Ausencia de perspectiva relacional.} Los an{\'a}lisis tabulares no detectan
    comunidades de proveedores con comportamientos coordinados, ni patrones de concentraci{\'o}n
    que son visibles {\'u}nicamente en el grafo de relaciones empresa-organismo \cite{Wachs2025ML}.
    En el dominio CPV~71, con 2.031 empresas y 394 organismos, la estructura del mercado
    s{\'o}lo es observable mediante an{\'a}lisis de grafos.

  \item \textbf{Invisibilidad del contexto documental.} Los pliegos, resoluciones y
    memorias justificativas ---que contienen se{\~n}ales textuales clave sobre posibles
    irregularidades--- permanecen inaccesibles para los sistemas de control basados
    exclusivamente en campos estructurados \cite{NLPLegalSurvey2024}.
\end{itemize}

\section{Eficacia del an{\'a}lisis de datos para la detecci{\'o}n de riesgo: evidencia
emp{\'i}rica (PI-4)}
\label{sec:pi4}

La literatura cient{\'i}fica de los {\'u}ltimos cinco a{\~n}os ofrece evidencias s{\'o}lidas y
replicadas sobre la eficacia del an{\'a}lisis de datos para la detecci{\'o}n de irregularidades
en contrataci{\'o}n p{\'u}blica. Se resumen a continuaci{\'o}n las contribuciones m{\'a}s relevantes:

\citeA{Coviello2023Italian} presentan el estudio emp{\'i}rico de mayor rigor en el {\'a}mbito
europeo: utilizando 4.000 procedimientos de contrataci{\'o}n italiana con 500 adjudicaciones
corruptas confirmadas judicialmente como \textit{ground truth}, demuestran que modelos de
Random Forest entrenados con indicadores de red flags ---licitaci{\'o}n {\'u}nica, concentraci{\'o}n
de adjudicaciones, desviaci{\'o}n de importe--- alcanzan precisiones del 70--75\,\% en
detecci{\'o}n de contratos corruptos. Su trabajo identifica, adem{\'a}s, los indicadores de
mayor poder predictivo individual y sus interacciones, proporcionando una base metodol{\'o}gica
directamente transferible al contexto espa{\~n}ol.

\citeA{Wachs2025ML} aplican un enfoque combinado de machine learning y an{\'a}lisis de redes
a datos de contratos p{\'u}blicos en M{\'e}xico, demostrando que las m{\'e}tricas de red
(centralidad de eigenvector, grado ponderado, coeficiente de agrupamiento) a{\~n}aden poder
predictivo significativo respecto a los modelos puramente tabulares. Su trabajo es
especialmente relevante porque utiliza un dise{\~n}o cuasiexperimental con proveedores
sancionados como etiquetas positivas, analogable a la situaci{\'o}n de datos disponibles
en Espa{\~n}a.

El proyecto europeo DIGIWHIST \cite{Fazekas2020DIGIWHIST,DIGIWHISTProject} desarroll{\'o}
el \textit{Corruption Risk Index} (CRI), combinando cuatro indicadores observables
---riesgo de licitaci{\'o}n, conexiones pol{\'i}ticas, riesgo del proveedor y riesgo del
{\'o}rgano contratante--- y lo valid{\'o} sobre datos de 35 jurisdicciones europeas. El CRI
est{\'a} disponible en OpenTender.eu \cite{OpenTenderEU,FazekasOpentender2025} para todos
los pa{\'i}ses de la UE, incluida Espa{\~n}a, y constituye el est{\'a}ndar de referencia para
indicadores de riesgo comparables.

El \textit{mapping study} sistem{\'a}tico de \citeA{FraudDetectionSystematic2025} (EPJ Data
Science, 2025) analiza m{\'a}s de 150 publicaciones sobre detecci{\'o}n de fraude en contrataci{\'o}n
p{\'u}blica y concluye que los modelos que combinan m{\'u}ltiples indicadores de red flags con
m{\'e}tricas de red superan consistentemente en rendimiento a los enfoques de indicador
{\'u}nico, con una mejora media del 15--20\,\% en AUC-ROC.

Finalmente, \citeA{MunozPla2026} demuestran que el \textit{dataset} BOE 2014--2024
presenta una calidad y coherencia temporal suficientes para el an{\'a}lisis de patrones
estructurales y temporales del riesgo en contrataci{\'o}n espa{\~n}ola, reforzando la
viabilidad t{\'e}cnica del proyecto.

\section{Dominio de an{\'a}lisis: CPV~71 como caso de validaci{\'o}n metodol{\'o}gicamente
justificado}
\label{sec:dominio_cpv71}

Tras el an{\'a}lisis exploratorio del \textit{dataset} BOE 2014--2024, se ha seleccionado la
\textbf{divisi{\'o}n CPV~71} (Servicios de arquitectura, construcci{\'o}n, ingenier{\'i}a e
inspecci{\'o}n) como dominio de validaci{\'o}n del prototipo ProcureWatch Analytics \cite{CPVCodes71}.
Esta divisi{\'o}n comprende 3.443 contratos adjudicados con un importe total superior a los
2.179~millones de euros, distribuidos entre 394 organismos contratantes y 2.031 empresas
adjudicatarias distintas.

La elecci{\'o}n no se fundamenta en el volumen absoluto del dominio, sino en tres criterios
metodol{\'o}gicos que lo convierten en el banco de pruebas {\'o}ptimo para evaluar los tres
m{\'o}dulos principales de la plataforma:

\begin{description}
  \item[Criterio~1: Estabilidad de los indicadores de riesgo tabulares.] A diferencia
    del CPV~45 (Obras, $n = 11.808$), los servicios de ingenier{\'i}a presentan una menor
    incidencia de modificaciones contractuales legalmente justificadas (art. 205 LCSP).
    En obras, las desviaciones de importe y plazo son fen{\'o}menos frecuentes y a menudo
    leg{\'i}timos por imprevistos t{\'e}cnicos o geopol{\'i}ticos, lo que introduce ruido en los
    modelos de detecci{\'o}n de anomal{\'i}as no supervisados y eleva la tasa de falsos positivos.
    Los servicios de ingenier{\'i}a, estructurados mayoritariamente en torno a tarifas horarias
    y entregables definidos, ofrecen una base m{\'a}s estable para los indicadores RF-03
    (desviaci{\'o}n de importe) y RF-04 (plazo de resoluci{\'o}n inusualmente corto)
    \cite{OCP2024RedFlags}. Esta propiedad es alineada con la metodolog{\'i}a del proyecto
    DIGIWHIST, que segmenta el an{\'a}lisis por mercados espec{\'i}ficos para garantizar la
    comparabilidad de los indicadores \cite{DIGIWHISTProject}.

  \item[Criterio~2: Riqueza del grafo para an{\'a}lisis de redes.] Con 2.031 empresas para
    3.443 contratos (ratio 1,7 contratos/empresa) y una concentraci{\'o}n del Top~1 inferior
    al 1\,\%, el mercado CPV~71 presenta la combinaci{\'o}n te{\'o}ricamente deseable: no tan
    concentrado como para que el grafo sea trivial, ni tan fragmentado como para que no
    existan patrones de comunidad detectables. Adem{\'a}s, el sector presenta una alta
    propensidad a la formaci{\'o}n de Uniones Temporales de Empresas (UTEs), que enriquecen
    el grafo bipartito empresa-organismo con relaciones de colaboraci{\'o}n recurrente entre
    competidores, precisamente las m{\'a}s relevantes para la detecci{\'o}n de colusi{\'o}n
    \cite{Wachs2025ML,Traag2019Leiden}.

  \item[Criterio~3: Adecuaci{\'o}n del corpus documental al m{\'o}dulo NLP.] Los pliegos de
    prescripciones t{\'e}cnicas en el {\'a}mbito de la ingenier{\'i}a contienen con frecuencia
    criterios de adjudicaci{\'o}n sujetos a juicio de valor ---experiencia en proyectos
    similares, metodolog{\'i}a propuesta, cualificaci{\'o}n del equipo--- que pueden representar
    hasta el 49\,\% de la puntuaci{\'o}n total en procedimientos abiertos con m{\'u}ltiples
    criterios. Estos criterios son extra{\'i}bles mediante NLP sobre los anuncios del BOE y
    los PDFs en PLACE, constituyendo el corpus id{\'o}neo para detectar posibles
    direccionamientos t{\'e}cnicos (\textit{tailored tenders}). La especificidad del vocabulario
    t{\'e}cnico (c{\'o}digos CPV de 8 d{\'i}gitos como 71311000 ---Consultor{\'i}a en ingenier{\'i}a
    civil--- o 71356000 ---Servicios t{\'e}cnicos---) permite, adem{\'a}s, una segmentaci{\'o}n
    precisa que reduce la varianza esp{\'u}ria en la comparaci{\'o}n de importes y la extracci{\'o}n
    de entidades \cite{CPVCodes71,NLPLegalSurvey2024}.
\end{description}

\section{Brecha identificada y posicionamiento del proyecto}

La Tabla~\ref{tab:gap_literatura} sintetiza el gap identificado en la literatura que
justifica el desarrollo de ProcureWatch Analytics.

\begin{table}[ht]
\centering
\small
\caption{Gap en la literatura y contribuci{\'o}n de ProcureWatch Analytics}
\label{tab:gap_literatura}
\begin{tabular}{p{3.2cm}p{4.5cm}p{4.5cm}}
\toprule
\textbf{Dimensi{\'o}n} & \textbf{Estado de la literatura} & \textbf{Contribuci{\'o}n de este TFM} \\
\midrule
Red flags en contrataci{\'o}n & Bien documentados y taxonomizados \cite{OCP2024RedFlags} & Implementaci{\'o}n, evaluaci{\'o}n y combinaci{\'o}n con ML \\
\addlinespace
ML para detecci{\'o}n de fraude & Demostrado emp{\'i}ricamente en Italia, M{\'e}xico, Kenia \cite{Coviello2023Italian,Wachs2025ML,NgangaHybrid2025} & Adaptaci{\'o}n y evaluaci{\'o}n sobre datos espa{\~n}oles \\
\addlinespace
An{\'a}lisis de redes & Probado en Portugal, Europa \cite{Martins2022Network,Fazekas2020DIGIWHIST} & Aplicaci{\'o}n sobre CPV~71 espa{\~n}ol con Leiden \\
\addlinespace
NLP en contrataci{\'o}n & Survey te{\'o}rico reciente, sin prototipo \cite{NLPLegalSurvey2024} & Implementaci{\'o}n sobre pliegos en espa{\~n}ol \\
\addlinespace
Plataforma integrada & No existe en datos espa{\~n}oles & Propuesta de este TFM \\
\addlinespace
Reproducibilidad & Limitada en estudios existentes & 100\,\% open-source + Docker \\
\addlinespace
Explicabilidad (AI Act) & Exigida desde 2024 \cite{AIAct2024} & SHAP + LIME con trazabilidad completa \\
\bottomrule
\end{tabular}
\end{table}

De este an{\'a}lisis se desprende que la brecha no reside en la ausencia de conocimiento
metodol{\'o}gico ---la literatura es rica y rigurosa--- sino en la ausencia de una
implementaci{\'o}n integrada, local, reproducible y evaluada sobre datos espa{\~n}oles.
ProcureWatch Analytics propone cerrar precisamente ese vac{\'i}o.

%=========================================================
\chapter{Contexto y Estado del Arte}
\label{chap:estadodelarte}
%=========================================================

\section{Organizaci{\'o}n de la revisi{\'o}n y criterios de inclusi{\'o}n}

La revisi{\'o}n del estado del arte cubre cinco bloques tem{\'a}ticos, directamente alineados
con los cinco componentes t{\'e}cnicos de ProcureWatch Analytics:
(A)~datos abiertos y est{\'a}ndar OCDS;
(B)~red flags, indicadores de riesgo y machine learning;
(C)~an{\'a}lisis de redes en contrataci{\'o}n p{\'u}blica;
(D)~NLP y procesamiento de documentos contractuales; y
(E)~visualizaci{\'o}n anal{\'i}tica y explicabilidad.

Se han priorizado publicaciones en revistas indexadas (Springer Nature, Cambridge
University Press, IEEE, ACM, Wiley) y proyectos financiados por la UE (DIGIWHIST,
OCP) del per{\'i}odo 2019--2026. Se incluyen \textit{preprints} de arXiv {\'u}nicamente
cuando el trabajo es ampliamente citado o representa la {\'u}nica fuente disponible sobre
un dato espec{\'i}fico (e.g., el dataset BOE 2014--2024 de \citeA{MunozPla2026}).

%---------------------------------------------------
\section{Bloque~A: Open Data de contrataci{\'o}n p{\'u}blica y est{\'a}ndar OCDS}
%---------------------------------------------------

\subsection{El est{\'a}ndar OCDS como eje de integraci{\'o}n}

El Open Contracting Data Standard (OCDS) \cite{OCDS2024} es el marco de referencia
internacional para la estructuraci{\'o}n y la interoperabilidad de datos de contrataci{\'o}n
p{\'u}blica. Define m{\'a}s de 40 campos por contrato ---organizados en cinco secciones:
\textit{planning}, \textit{tender}, \textit{award}, \textit{contract} y
\textit{implementation}--- y ha sido adoptado por m{\'a}s de 60 pa{\'i}ses. El Banco
Mundial \cite{WorldBankOCDS2021} ha publicado gu{\'i}as de implementaci{\'o}n que detallan
c{\'o}mo armonizar fuentes nacionales heterog{\'e}neas con el est{\'a}ndar; sus directrices sobre
identificadores de org{\'a}nismo y de empresa son directamente aplicables al contexto de
PLACE y el BOE.

El an{\'a}lisis comparativo de la OCP sobre apertura de datos en la UE
\cite{OCP2023EUOpenData} constata que Espa{\~n}a ha avanzado notablemente en la
publicaci{\'o}n de datos, pero subsisten problemas de completitud en campos cr{\'i}ticos
(NIF del adjudicatario, importe adjudicado exacto, fechas de hitos) y de inconsistencias
entre la informaci{\'o}n publicada a nivel de expediente y a nivel de lote. Estas
deficiencias definen el punto de partida y las tareas de limpieza del Agente~1
(\textit{pipeline} de datos) de ProcureWatch.

\subsection{El dataset BOE 2014--2024 como fuente longitudinal}

\citeA{MunozPla2026} presentan el primer \textit{dataset} longitudinal de contrataci{\'o}n
p{\'u}blica espa{\~n}ola en formato OCDS: aproximadamente 97.000 contratos adjudicados
publicados en el BOE entre 2014 y 2024, con informaci{\'o}n sobre organismos contratantes,
empresas adjudicatarias, importes, procedimientos y c{\'o}digos CPV. La publicaci{\'o}n del
\textit{dataset} como recurso abierto y reproducible es una contribuci{\'o}n metodol{\'o}gica
relevante: permite replicar an{\'a}lisis, comparar resultados entre estudios y evaluar la
evoluci{\'o}n temporal del sistema de contrataci{\'o}n a lo largo de una d{\'e}cada.

Para el presente trabajo, el dataset BOE 2014--2024 constituye la fuente de datos
primaria del prototipo, complementada con datos en tiempo real de PLACE.

\subsection{Plataformas europeas de referencia: OpenTender y TED}

A nivel europeo, la plataforma OpenTender.eu \cite{OpenTenderEU,FazekasOpentender2025}
centraliza datos de 35 jurisdicciones con m{\'a}s de 19 millones de contratos e indicadores
de riesgo precalculados, incluido el \textit{Corruption Risk Index} (CRI) desarrollado
por el proyecto DIGIWHIST \cite{Fazekas2020DIGIWHIST}. Para Espa{\~n}a, incluye datos
desde 2007 en formato JSON con campos OCDS. El portal TED (\textit{Tenders Electronic
Daily}) recoge contratos que superan los umbrales comunitarios de publicidad, con
cobertura desde 1994. Estas plataformas complementan el dataset BOE y permiten
comparativas de riesgo en contexto europeo, lo que enriquece la validaci{\'o}n del
prototipo.

%---------------------------------------------------
\section{Bloque~B: Red flags, indicadores de riesgo y machine learning}
%---------------------------------------------------

\subsection{Taxonom{\'i}a de red flags: estado actual}

La Open Contracting Partnership \cite{OCP2024RedFlags} ha publicado la gu{\'i}a de
referencia m{\'a}s actualizada sobre indicadores de riesgo en contrataci{\'o}n p{\'u}blica.
Define m{\'a}s de 150 red flags organizados por fase del ciclo contractual (planificaci{\'o}n,
licitaci{\'o}n, adjudicaci{\'o}n, ejecuci{\'o}n), tipo de irregularidad (fraccionamiento,
colusi{\'o}n, conflicto de inter{\'e}s, modificaci{\'o}n injustificada) y tipo de datos necesarios
para su c{\'a}lculo (campos OCDS, datos de empresa, datos relacionales). La taxonom{\'i}a
original de 2016 \cite{OCP2016RedFlags} ha sido revisada y ampliada en 2024 incorporando
la dimensi{\'o}n de redes societarias y el contexto de datos masivos.

Los indicadores m{\'a}s frecuentemente implementados en la literatura, y que constituyen
el n{\'u}cleo del motor de red flags de ProcureWatch, son:
\begin{enumerate}
  \item Tasa de licitaci{\'o}n {\'u}nica o restringida (RF-01).
  \item Concentraci{\'o}n relativa de adjudicaciones por organismo (RF-02).
  \item Desviaci{\'o}n porcentual entre importe estimado y adjudicado (RF-03).
  \item Plazo de resoluci{\'o}n inusualmente corto (RF-04).
  \item Recurrencia del mismo proveedor en procedimientos negociados (RF-05).
  \item Uso reiterado del contrato menor en importes cercanos al umbral (RF-06).
\end{enumerate}

\subsection{Machine learning: de los indicadores aislados a los modelos compuestos}

El estudio de \citeA{Coviello2023Italian} es el punto de referencia metodol{\'o}gico m{\'a}s
s{\'o}lido disponible para el contexto europeo. Utilizando datos de licitaciones italianas
con 500 contratos corruptos confirmados judicialmente, demuestran que un Random Forest
entrenado sobre los indicadores de la taxonom{\'i}a OCP alcanza el 70--75\,\% de precisi{\'o}n,
frente a un 45--50\,\% de los indicadores aislados. Su an{\'a}lisis de importancia de
variables mediante Gini revela que la licitaci{\'o}n {\'u}nica y la concentraci{\'o}n de
adjudicaciones son los predictores individuales m{\'a}s potentes, pero que sus interacciones
a{\~n}aden hasta un 12\,\% de mejora adicional en AUC-ROC.

\citeA{NgangaHybrid2025} proponen un \textit{framework} h{\'i}brido para datos de
contrataci{\'o}n en Kenia combinando Logistic Regression, Random Forest y \textit{clustering}
no supervisado, obteniendo mejoras sobre cada modelo individual. \citeA{DataDrivenTransparency2025}
incorporan an{\'a}lisis de redes sociales como capa adicional de \textit{features}, con
mejoras adicionales en AUC cuando se a{\~n}aden m{\'e}tricas de centralidad. El
\textit{mapping study} sistem{\'a}tico de \citeA{FraudDetectionSystematic2025} confirma,
sobre m{\'a}s de 150 publicaciones, que los modelos compuestos superan consistentemente a
los enfoques de indicador {\'u}nico.

El framework PANG de \citeA{LiuPatternMining2023} es espec{\'i}ficamente relevante para
ProcureWatch: aplica \textit{pattern mining} sobre grafos de contrataci{\'o}n para detectar
subgrafos an{\'o}malos ---conjuntos de contratos, organismos y empresas cuya estructura de
relaciones es estad{\'i}sticamente improbable en un mercado competitivo--- logrando detectar
tramas de adjudicaciones coordinadas no identificables con indicadores tabulares.

\subsection{Positive-Unlabeled Learning: tratamiento de la ausencia de etiquetas negativas}

Un desaf{\'i}o metodol{\'o}gico central en la detecci{\'o}n de fraude contractual es la ausencia
de etiquetas negativas claras. Se dispone de algunos casos confirmados de irregularidad
(etiquetas positivas) y de un amplio volumen de contratos sin etiqueta, pero no es
l{\'i}cito asumir que todos los contratos no etiquetados son leg{\'i}timos. El enfoque
\textit{Positive-Unlabeled Learning} (PUL) \cite{PULearningIJCAI2021} aborda este
problema entrenando clasificadores que aprendan a distinguir positivos de instancias no
etiquetadas, sin penalizar como errores los positivos no identificados. \citeA{FraudPULearningIEEE2023}
demuestran su eficacia con un enfoque contrastivo aplicado espec{\'i}ficamente a detecci{\'o}n
de fraude financiero, con mejoras de hasta 8 puntos porcentuales en precisi{\'o}n sobre
m{\'e}todos supervisados cl{\'a}sicos cuando la tasa de casos etiquetados es inferior al 5\,\%.

\subsection{Detecci{\'o}n de anomal{\'i}as no supervisada}

Para el escenario en que no existen etiquetas positivas conocidas, m{\'e}todos no
supervisados como Isolation Forest \cite{IsolationForestSklearn} permiten identificar
contratos at{\'i}picos en funci{\'o}n de la distancia de sus caracter{\'i}sticas al resto del
conjunto. \citeA{OcampoGutierrez2019} aplican este enfoque directamente sobre datos
OCDS de contrataci{\'o}n paraguaya, obteniendo resultados coherentes con el conocimiento
experto sobre irregularidades conocidas. En ProcureWatch, Isolation Forest se utiliza
como l{\'i}nea de base para comparar con los enfoques supervisados y PUL.

%---------------------------------------------------
\section{Bloque~C: An{\'a}lisis de redes en contrataci{\'o}n p{\'u}blica}
%---------------------------------------------------

\subsection{Grafos bipartitos empresa-organismo y m{\'e}tricas de centralidad}

El grafo bipartito empresa-organismo-contrato es la estructura topol{\'o}gica natural para
representar el sistema de adjudicaciones: los nodos son empresas y organismos; las
aristas son contratos adjudicados, con atributos de importe, tipo de procedimiento y
fecha. Esta representaci{\'o}n permite calcular m{\'e}tricas de centralidad que capturan
propiedades del mercado no observables en tablas planas:
\begin{itemize}
  \item \textbf{Grado ponderado por importe}: suma de los importes adjudicados por nodo,
    indicadora de participaci{\'o}n econ{\'o}mica en el mercado.
  \item \textbf{Centralidad de intermediaci{\'o}n (\textit{betweenness})}: empresas que
    actúan como ``puentes'' en la red, adjudicatarias de m{\'u}ltiples organismos de
    perfil heterog{\'e}neo, lo que puede indicar capacidad de influencia transversal.
  \item \textbf{Centralidad de vector propio (\textit{eigenvector})}: empresas conectadas
    a organismos con alta actividad contractual, relevante para detectar relaciones
    preferenciales de alto valor.
\end{itemize}

\citeA{Wachs2025ML} demuestran que estas m{\'e}tricas, combinadas con indicadores de red
flags, mejoran la predicci{\'o}n de riesgo de forma significativa y estable en diferentes
contextos nacionales. \citeA{Martins2022Network} aplican an{\'a}lisis de grafos a la
contrataci{\'o}n portuguesa, identificando comunidades de riesgo cuya detecci{\'o}n no es
posible mediante indicadores tabulares.

\citeA{GraphSIF2020} aportan evidencia desde el contexto de redes B2B: demuestran que
el an{\'a}lisis de flujos de pago en grafos bipartitos permite detectar suplantaci{\'o}n de
proveedores con alta precisi{\'o}n. \citeA{GroupFraudKDD2023} extienden este enfoque a
la detecci{\'o}n de fraude colectivo en grafos bipartitos de plataformas de comercio
electr{\'o}nico, con m{\'e}todos directamente transferibles al contexto de contrataci{\'o}n.

\subsection{Detecci{\'o}n de comunidades: del algoritmo Louvain al Leiden}

La detecci{\'o}n de comunidades en el grafo empresa-organismo permite identificar grupos de
empresas que licitan sistem{\'a}ticamente en los mismos organismos ---posible indicador de
colusi{\'o}n o acuerdo de reparto de mercado--- o grupos de organismos que adjudican
recurrentemente a las mismas empresas. El indicador est{\'a}ndar de calidad de la
partici{\'o}n comunitaria es la modularidad~$Q$: valores superiores a 0,30 se consideran
indicativos de estructura comunitaria real, no artefactual \cite{Traag2019Leiden}.

\citeA{Traag2019Leiden} presentan el algoritmo Leiden como mejora demostrada del
Louvain: garantiza comunidades internas bien conectadas (resolviendo un fallo estructural
del Louvain donde comunidades aparentes pod{\'i}an estar internamente desconectadas) y
reduce el tiempo de convergencia en grafos grandes. La implementaci{\'o}n eficiente en
memoria compartida de \citeA{FastLeiden2024} hace viable la ejecuci{\'o}n del Leiden sobre
el grafo completo de contrataci{\'o}n espa{\~n}ola sin necesidad de infraestructura distribuida.

\subsection{Relaciones societarias: UTEs como enriquecimiento del grafo}

En el dominio CPV~71, la alta propensidad a la formaci{\'o}n de Uniones Temporales de
Empresas (UTEs) introduce un tercer tipo de nodo en el grafo: la entidad temporal
formada por empresas que compiten de forma independiente en otros procedimientos.
La detecci{\'o}n de empresas que forman UTEs recurrentes con la misma composici{\'o}n, o que
forman UTEs con empresas con las que tienen relaciones societarias previas, es un
indicador de posible coordinaci{\'o}n que complementa los indicadores tabulares. Este
enriquecimiento del grafo es una aportaci{\'o}n metodol{\'o}gica espec{\'i}fica del presente
trabajo.

%---------------------------------------------------
\section{Bloque~D: NLP y procesamiento de documentos de contrataci{\'o}n}
%---------------------------------------------------

\subsection{El dominio legal como reto espec{\'i}fico para el NLP}

El procesamiento de documentos de contrataci{\'o}n ---pliegos de prescripciones t{\'e}cnicas,
resoluciones de adjudicaci{\'o}n, memorias justificativas--- plantea retos espec{\'i}ficos
respecto al NLP de prop{\'o}sito general. El survey de \citeA{NLPLegalSurvey2024}, que
revisa 3.000 publicaciones sobre NLP legal publicadas hasta 2024, identifica cinco
dimensiones que hacen del dominio legal un entorno singular: (1)~vocabulario especializado
y jerga jur{\'i}dica con significados distintos a los del lenguaje general; (2)~referencias
cruzadas a art{\'i}culos legales, procedimientos y entidades que exigen resoluci{\'o}n de
co-referencia; (3)~documentos largos (10--200 p{\'a}ginas) con estructura jerarquizada;
(4)~necesidad de extraer fragmentos espec{\'i}ficos en posiciones predecibles (art{\'i}culos,
cl{\'a}usulas); y (5)~baja disponibilidad de corpus etiquetados para entrenamiento.

En el contexto de contrataci{\'o}n p{\'u}blica en espa{\~n}ol, estos retos se ven amplificados
por la variabilidad de formatos entre organismos y la mezcla de texto leg{\'i}ble por
m{\'a}quina (pliegos en Word o HTML) y PDFs escaneados que requieren OCR previo.

\subsection{Reconocimiento de entidades nombradas (NER) en documentos administrativos}

La tarea central del m{\'o}dulo NLP de ProcureWatch es el reconocimiento de entidades
nombradas (NER): identificar y clasificar en texto no estructurado los organismos
contratantes, empresas licitadoras, importes, plazos, criterios de adjudicaci{\'o}n y
c{\'o}digos CPV. \citeA{NERLegalDocuments2020} proponen enfoques de granularidad fina
para NER en documentos legales, con especial atenci{\'o}n a la fragmentaci{\'o}n de
entidades largas y a la co-referencia entre nombres completos y abreviaciones.

Para el espa{\~n}ol, el framework NLP de referencia en producci{\'o}n es spaCy \cite{SpaCyOfficial},
con soporte nativo para espa{\~n}ol a trav{\'e}s del modelo \texttt{es\_core\_news\_lg}. Sus
capacidades incluyen tokenizaci{\'o}n, POS tagging, an{\'a}lisis de dependencias y NER sobre
18 tipos de entidades. Para el dominio espec{\'i}fico de contrataci{\'o}n, se requiere
\textit{fine-tuning} adicional sobre un corpus del dominio para alcanzar las precisiones
reportadas en la literatura ($>75$\,\% en F1 para entidades cr{\'i}ticas). Los modelos de
transformers (BERT-base, Legal-BETO) aportan comprensi{\'o}n contextual que mejora el NER
en casos ambig{\"u}os, a costa de mayor coste computacional \cite{NLPLegalSurvey2024}.

\subsection{OCR y parseado de documentos administrativos}

Una parte sustancial de los documentos disponibles en PLACE est{\'a} en formato PDF
escaneado, lo que requiere OCR previo a cualquier an{\'a}lisis NLP. Docling \cite{DoclingIBM2024}
(IBM Research, 2024) es la herramienta open-source de referencia para parsing avanzado
de documentos complejos con OCR integrado: maneja tablas, encabezados, notas al pie y
layouts mult{\'i}columna, y exporta en formatos JSON y Markdown directamente compatibles
con pipelines de NLP. El benchmark comparativo de 2025 \cite{DoclingIBM2024} posiciona
a Docling como la alternativa open-source con mejor relaci{\'o}n precisi{\'o}n/velocidad para
documentos de tipo administrativo o legal, frente a alternativas cloud-based como
Amazon Textract.

\subsection{Modelos de lenguaje locales para an{\'a}lisis jur{\'i}dico-administrativo}

Los modelos de lenguaje de gran escala (LLMs) en ejecuci{\'o}n local permiten analizar
documentos sin dependencia de APIs externas, preservando la confidencialidad de los
datos y facilitando el cumplimiento con el RGPD y el AI Act. Para el espa{\~n}ol, Qwen3-7B
\cite{AlibabaQwen3} ofrece un soporte de alta calidad con requerimientos de memoria
gestionables (64~GB RAM); Mistral Small~3 \cite{MistralSmall3} es la referencia
comparativa para tareas de s{\'i}ntesis y clasificaci{\'o}n. En ProcureWatch, el LLM local
tiene un rol estrictamente auxiliar y acotado: estructurar documentos, extraer
entidades con contexto y generar explicaciones de casos. La decisi{\'o}n de si un contrato
es sospechoso es siempre humana, en cumplimiento del art{\'i}culo 14 del AI Act \cite{AIAct2024}.

\subsection{Recuperaci{\'o}n aumentada con generaci{\'o}n (RAG)}

El sistema RAG de ProcureWatch combina un recuperador de fragmentos documentales
sem{\'a}nticamente relevantes con el LLM local para producir respuestas contextualizadas
sobre un caso de riesgo: dado un contrato o adjudicatario con m{\'u}ltiples red flags
activados, el motor recupera y sintetiza los fragmentos de pliegos y resoluciones que
explican o matizan los indicadores activados \cite{LangChain2024LangGraph}.

La recuperaci{\'o}n h{\'i}brida combina dos enfoques complementarios: b{\'u}squeda densa
mediante embeddings multilingu{\"e}s BGE-M3 \cite{BGEAM3Embeddings} (indexados en Qdrant)
y b{\'u}squeda l{\'e}xica BM25 para t{\'e}rminos t{\'e}cnicos espec{\'i}ficos de contrataci{\'o}n
(como c{\'o}digos CPV o referencias a art{\'i}culos de la LCSP) que no son capturados
eficientemente por la b{\'u}squeda sem{\'a}ntica. La calidad del sistema se eval{\'u}a con el
framework RAGAS, con objetivos: Faithfulness $>0{,}80$, Answer Relevancy $>0{,}75$ y
Context Recall $>0{,}70$.

%---------------------------------------------------
\section{Bloque~E: Visualizaci{\'o}n anal{\'i}tica y explicabilidad}
%---------------------------------------------------

\subsection{Explicabilidad del scoring de riesgo: requisito legal y t{\'e}cnico}

El Reglamento de IA de la UE \cite{AIAct2024} clasifica como sistemas de \textbf{alto
riesgo} aquellos que generan puntuaciones que puedan influir en decisiones sobre
empresas o contratos p{\'u}blicos (Anexo III, secci{\'o}n 8). Esta clasificaci{\'o}n implica
obligaciones espec{\'i}ficas: trazabilidad de las decisiones, documentaci{\'o}n del sistema,
evaluaci{\'o}n de conformidad y, fundamentalmente, supervisi{\'o}n humana obligatoria. En
consecuencia, la explicabilidad no es una caracter{\'i}stica opcional del prototipo, sino
un requisito normativo.

Los m{\'e}todos SHAP (\textit{SHapley Additive exPlanations}) y LIME (\textit{Local
Interpretable Model-Agnostic Explanations}) \cite{SHAPLIMEWiley2025} son los enfoques
de referencia para la explicabilidad de modelos de machine learning. SHAP calcula la
contribuci{\'o}n marginal exacta de cada \textit{feature} al score de riesgo de un contrato
espec{\'i}fico, basada en la teor{\'i}a de valores de Shapley; LIME construye una aproximaci{\'o}n
local lineal del modelo en el entorno del caso de inter{\'e}s. En ProcureWatch, SHAP se
utiliza para la ficha de riesgo autom{\'a}tica y para el panel de contribuci{\'o}n de
indicadores; LIME como validaci{\'o}n de coherencia de las explicaciones SHAP.

\subsection{Visualizaci{\'o}n del grafo empresa-organismo en \textit{front-end} web}

Para la visualizaci{\'o}n interactiva del grafo bipartito en el navegador, el benchmarking
de 2026 \cite{GraphVizPkgPulse2026} posiciona a Sigma.js como la mejor opci{\'o}n para
grafos grandes ($>10.000$ nodos) gracias a su motor WebGL de renderizado en GPU, y a
Cytoscape.js como alternativa para grafos medianos ($<5.000$ nodos) con mayores
necesidades de an{\'a}lisis interactivo integrado. En ProcureWatch, Sigma.js se utiliza
para el grafo principal y PyVis (\textit{vis-network} en Python) para la visualizaci{\'o}n
r{\'a}pida integrada en el panel Streamlit durante el prototipado.

\subsection{Dise{\~n}o de paneles de riesgo}

El dise{\~n}o del panel de riesgo de ProcureWatch sigue los principios de jerarqu{\'i}a de
informaci{\'o}n y flujo de revisi{\'o}n documentados en la literatura de \textit{procurement
analytics} \cite{IvaluaDashboard}: (1)~priorizar el ranking de casos por nivel de riesgo
en la vista principal; (2)~facilitar la navegaci{\'o}n desde el caso general al detalle del
contrato mediante selecciones en el grafo; (3)~incorporar la explicaci{\'o}n de los
indicadores activados en la ficha del caso, sin requerir conocimiento t{\'e}cnico del
modelo. El objetivo de usabilidad SUS\,$>68$ corresponde al umbral ``aceptable'' de la
escala de Brooke, validado sobre usuarios con perfil de analista t{\'e}cnico no experto en
ML.

%---------------------------------------------------
\section{Conclusiones del estado del arte: cuatro hallazgos que justifican el proyecto}
%---------------------------------------------------

La revisi{\'o}n sistem{\'a}tica de la literatura permite extraer cuatro hallazgos que justifican
directamente la propuesta:

\begin{enumerate}
  \item \textbf{Datos disponibles de alta calidad, sin explotaci{\'o}n sistem{\'a}tica.}
    El \textit{dataset} BOE 2014--2024 \cite{MunozPla2026} y los datos abiertos de PLACE
    \cite{MinHacienda2024PLACE} constituyen una base de datos de contrataci{\'o}n de alta
    calidad y cobertura decenal. Sin embargo, no existe a la fecha ninguna herramienta
    open-source que los integre y analice con un enfoque combinado de red flags, an{\'a}lisis
    de redes y NLP.

  \item \textbf{Eficacia demostrada de los m{\'e}todos combinados.} La literatura demuestra
    consistentemente, en tres contextos europeos y latinoamericanos distintos
    \cite{Coviello2023Italian,Wachs2025ML,Fazekas2020DIGIWHIST}, que los modelos que
    combinan indicadores de red flags con m{\'e}tricas de red mejoran la detecci{\'o}n de
    irregularidades en 15--20 puntos de AUC-ROC respecto a indicadores aislados. Esta
    mejora no se ha validado sobre datos espa{\~n}oles.

  \item \textbf{NLP como dimensi{\'o}n infraexplotada.} A pesar de que el survey de
    \citeA{NLPLegalSurvey2024} documenta una madurez creciente de las t{\'e}cnicas, el
    procesamiento de documentos contractuales (pliegos, resoluciones) con NLP es
    pr{\'a}cticamente inexistente en la literatura de detecci{\'o}n de fraude en contrataci{\'o}n
    p{\'u}blica, a pesar de que, como se ha documentado en la Secci{\'o}n~\ref{sec:pi2},
    contienen informaci{\'o}n cualitativa sobre el riesgo que no es capturada por los datos
    estructurados.

  \item \textbf{Ausencia de prototipo integrado, local y evaluado.} No existe en la
    literatura ning{\'u}n prototipo open-source que combine \textit{pipeline} OCDS, motor
    de red flags con ML, an{\'a}lisis de redes con Leiden, m{\'o}dulo NLP con RAG y front-end
    explicable sobre datos de contrataci{\'o}n p{\'u}blica espa{\~n}ola. ProcureWatch Analytics
    propone cubrir exactamente esta brecha.
\end{enumerate}

%=========================================================
\chapter{Objetivos}
%=========================================================

\section{Objetivo general}

Dise{\~n}ar, desarrollar y evaluar \textbf{ProcureWatch Analytics}, una plataforma
anal{\'i}tica multiagente que, a partir de datos abiertos de contrataci{\'o}n p{\'u}blica
espa{\~n}ola (BOE 2014--2026 y PLACE) y del est{\'a}ndar OCDS, detecte, cuantifique y
visualice patrones de riesgo e indicadores de red flags en licitaciones, adjudicaciones
y empresas adjudicatarias del dominio CPV~71, mediante an{\'a}lisis de datos masivos,
an{\'a}lisis de redes societarias, procesamiento de lenguaje natural sobre documentos
contractuales y modelos de scoring de anomal{\'i}as con explicabilidad verificable,
generando fichas de riesgo interpretables que apoyen la priorizaci{\'o}n de casos para
revisi{\'o}n por organismos de control, en cumplimiento del marco normativo vigente.

\section{Objetivos espec{\'i}ficos}

\subsection*{OE1 --- An{\'a}lisis del dominio y las fuentes de datos}

\textbf{Descripci{\'o}n:} Estudiar exhaustivamente el ciclo de contrataci{\'o}n p{\'u}blica
espa{\~n}ola, la estructura y calidad de los datos disponibles en PLACE y en el
\textit{dataset} BOE 2014--2024 seg{\'u}n est{\'a}ndar OCDS, con especial foco en el
dominio CPV~71, e identificar las variables m{\'a}s relevantes para la detecci{\'o}n de
red flags \cite{MunozPla2026,OCP2024RedFlags}.

\textbf{M{\'e}tricas de {\'e}xito verificables:}
\begin{itemize}
  \item Documento de an{\'a}lisis exploratorio que acredita completitud, coherencia temporal
    y distribuci{\'o}n de variables OCDS en el dominio CPV~71.
  \item Matriz de relevancia: mapeo de campos OCDS $\rightarrow$ indicadores de red flags
    RF-01 a RF-15, con justificaci{\'o}n bibliogr{\'a}fica de cada asignaci{\'o}n.
  \item M{\'i}nimo 3 fuentes de datos complementarias identificadas y documentadas
    (ROLECE, registros mercantiles, TED).
\end{itemize}

\subsection*{OE2 --- Dise{\~n}o del modelo de datos anal{\'i}tico}

\textbf{Descripci{\'o}n:} Definir e implementar un esquema de datos integrado y armonizado
con el est{\'a}ndar OCDS \cite{OCDS2024} que unifique informaci{\'o}n de contratos,
licitaciones, adjudicatarios, {\'o}rganos contratantes y relaciones societarias, como
fundamento del an{\'a}lisis de redes y del motor de red flags.

\textbf{M{\'e}tricas de {\'e}xito verificables:}
\begin{itemize}
  \item Modelo Entidad-Relaci{\'o}n documentado en PostgreSQL con m{\'i}nimo 6 entidades
    principales y claves externas verificadas.
  \item Grafo empresa-organismo-contrato construido en Neo4j con coherencia estructural
    demostrable.
  \item Completitud de campos OCDS cr{\'i}ticos $\geq 85$\,\% tras el proceso de normalizaci{\'o}n.
\end{itemize}

\subsection*{OE3 --- Construcci{\'o}n del \textit{pipeline} de datos}

\textbf{Descripci{\'o}n:} Implementar un \textit{pipeline} reproducible de descarga,
limpieza, normalizaci{\'o}n, enriquecimiento y carga de datos OCDS, con validaci{\'o}n
autom{\'a}tica de calidad y trazabilidad de fuentes, que garantice la reproducibilidad
integral del proceso \cite{WorldBankOCDS2021}.

\textbf{M{\'e}tricas de {\'e}xito verificables:}
\begin{itemize}
  \item \textit{Pipeline} ejecutable y documentado (Python\,+\,Pandas\,+\,DuckDB) con
    tests autom{\'a}ticos de calidad.
  \item Completitud OCDS post-carga $>90$\,\% en campos cr{\'i}ticos; coherencia temporal
    $>98$\,\%.
  \item Tiempo de ejecuci{\'o}n completo $<2$ horas para el \textit{dataset} BOE 2014--2024.
\end{itemize}

\subsection*{OE4 --- Desarrollo del motor de red flags}

\textbf{Descripci{\'o}n:} Dise{\~n}ar e implementar un conjunto de m{\'i}nimo 15 indicadores
de red flags (RF-01 a RF-15) basados en la taxonom{\'i}a de la OCP \cite{OCP2024RedFlags}
y en \citeA{Fazekas2020DIGIWHIST}, complementados con tres enfoques de machine learning
para scoring de riesgo compuesto: reglas expl{\'i}citas, Isolation Forest y
Positive-Unlabeled Learning.

\textbf{M{\'e}tricas de {\'e}xito verificables:}
\begin{itemize}
  \item M{\'i}nimo 15 indicadores documentados con umbral, l{\'o}gica de c{\'a}lculo y
    fundamento bibliogr{\'a}fico.
  \item Modelo de scoring compuesto que combina flags individuales con m{\'e}tricas de red.
  \item Precisi{\'o}n $>70$\,\% en un conjunto de casos con \textit{ground truth} parcial
    (contratos con resoluciones judiciales o denuncias p{\'u}blicas documentadas).
  \item Comparativa de los tres enfoques ML con m{\'e}tricas de AUC-ROC, precisi{\'o}n y
    estabilidad en validaci{\'o}n cruzada.
\end{itemize}

\subsection*{OE5 --- M{\'o}dulo de an{\'a}lisis de redes}

\textbf{Descripci{\'o}n:} Construir y analizar el grafo bipartito empresa-organismo-contrato
para detectar comunidades (algoritmo Leiden \cite{Traag2019Leiden}), patrones de
centralidad at{\'i}picos y posibles relaciones societarias de riesgo en el dominio CPV~71,
correlacionando las m{\'e}tricas de red con los scores del motor de red flags.

\textbf{M{\'e}tricas de {\'e}xito verificables:}
\begin{itemize}
  \item Grafo construido con m{\'i}nimo 2.031 nodos empresa y 394 organismos del dominio CPV~71.
  \item Ejecuci{\'o}n del algoritmo Leiden con modularidad $Q>0{,}30$.
  \item Identificaci{\'o}n documentada de m{\'i}nimo 5 patrones de concentraci{\'o}n at{\'i}pica.
  \item Correlaci{\'o}n estad{\'i}sticamente significativa ($\rho \geq 0{,}45$, $p < 0{,}05$)
    entre score de centralidad de red y score de red flags por empresa.
\end{itemize}

\subsection*{OE6 --- M{\'o}dulo NLP y recuperaci{\'o}n documental}

\textbf{Descripci{\'o}n:} Implementar un sistema de extracci{\'o}n de informaci{\'o}n
estructurada sobre pliegos de prescripciones t{\'e}cnicas y resoluciones de adjudicaci{\'o}n
del dominio CPV~71 mediante OCR (Docling \cite{DoclingIBM2024}), NER en espa{\~n}ol
(spaCy \cite{SpaCyOfficial}) y motor RAG h{\'i}brido (Qdrant + BM25 + BGE-M3
\cite{BGEAM3Embeddings}).

\textbf{M{\'e}tricas de {\'e}xito verificables:}
\begin{itemize}
  \item Corpus de m{\'i}nimo 5.000 documentos procesados con Docling y almacenados en
    el \textit{vector store}.
  \item NER en espa{\~n}ol con precisi{\'o}n $>75$\,\% en entidades cr{\'i}ticas (organismos,
    empresas, importes, c{\'o}digos CPV).
  \item Motor RAG con Faithfulness $>0{,}80$ y Answer Relevancy $>0{,}75$ (RAGAS).
  \item Tiempo de recuperaci{\'o}n de contexto por caso $<5$ segundos.
\end{itemize}

\subsection*{OE7 --- Desarrollo del \textit{front-end} anal{\'i}tico}

\textbf{Descripci{\'o}n:} Construir una interfaz web interactiva (Streamlit + Sigma.js +
Plotly) con grafo de relaciones, panel de indicadores de riesgo por caso, \textit{timeline}
de evoluci{\'o}n temporal y ficha de riesgo explicable generada autom{\'a}ticamente con SHAP
\cite{SHAPLIMEWiley2025} y contexto documental recuperado por el motor RAG.

\textbf{M{\'e}tricas de {\'e}xito verificables:}
\begin{itemize}
  \item Puntuaci{\'o}n SUS (\textit{System Usability Scale}) $>68$ en evaluaci{\'o}n con
    m{\'i}nimo 5 usuarios de perfil analista t{\'e}cnico.
  \item Tiempo medio para localizar y documentar un caso de riesgo $<5$ minutos.
  \item Ficha explicativa evaluada con claridad $\geq 3{,}5/5$ (escala Likert) por los
    mismos evaluadores.
\end{itemize}

\subsection*{OE8 --- Evaluaci{\'o}n integral del sistema}

\textbf{Descripci{\'o}n:} Validar el prototipo mediante m{\'e}tricas cuantitativas de calidad
del \textit{pipeline} (completitud OCDS), precisi{\'o}n del motor de red flags,
modularidad del grafo, m{\'e}tricas RAG y usabilidad, y an{\'a}lisis cualitativo de m{\'i}nimo
10 casos de riesgo en el dominio CPV~71 con \textit{ground truth} verificable.

\textbf{M{\'e}tricas de {\'e}xito verificables:}
\begin{itemize}
  \item Informe de evaluaci{\'o}n con todas las m{\'e}tricas especificadas en OE1--OE7.
  \item Fichas documentadas de m{\'i}nimo 10 casos de riesgo analizados, con descripci{\'o}n
    del patr{\'o}n detectado, indicadores activados y valoraci{\'o}n de la explicaci{\'o}n generada.
  \item Documento de limitaciones identificadas y \textit{roadmap} de mejora para
    trabajo futuro.
\end{itemize}

%=========================================================
\chapter{Marco Normativo}
%=========================================================

El dise{\~n}o de ProcureWatch Analytics se enmarca en un conjunto de disposiciones
normativas espa{\~n}olas y europeas que condicionan tanto el tratamiento de los datos como
la arquitectura del sistema de scoring:

\begin{description}
  \item[Ley 9/2017 de Contratos del Sector P{\'u}blico (LCSP)] \cite{LCSP2017}
    Marco regulatorio fundamental de la contrataci{\'o}n p{\'u}blica espa{\~n}ola. Define tipos
    de contrato, procedimientos, umbrales de publicidad obligatoria y los l{\'i}mites del
    contrato menor. Su art{\'i}culo 205 regula las modificaciones contractuales justificadas,
    clave para la calibraci{\'o}n del indicador RF-03 (desviaci{\'o}n de importe) en sectores
    donde las modificaciones son frecuentes por causas leg{\'i}timas (e.g., obras frente a
    servicios de ingenier{\'i}a).

  \item[Directiva 2014/24/UE del Parlamento Europeo y del Consejo] \cite{Directive2014}
    Marco europeo de contrataci{\'o}n p{\'u}blica que establece los principios comunitarios de
    concurrencia, transparencia, igualdad de trato y proporcionalidad. Define los umbrales
    de publicidad obligatoria en el DOUE (Diario Oficial de la UE), por encima de los
    cuales los contratos se publican tambi{\'e}n en TED.

  \item[Ley 19/2013 de transparencia, acceso a la informaci{\'o}n p{\'u}blica y buen gobierno]
    \cite{Ley192013}
    Establece la obligaci{\'o}n de publicar activamente la informaci{\'o}n sobre contratos del
    sector p{\'u}blico y define las condiciones de reutilizaci{\'o}n. Fundamenta jur{\'i}dicamente
    el uso de datos publicados en PLACE y el BOE para an{\'a}lisis cient{\'i}ficos y de
    supervisi{\'o}n.

  \item[RGPD y LOPDGDD]
    El Reglamento General de Protecci{\'o}n de Datos (RGPD) y la Ley Org{\'a}nica de
    Protecci{\'o}n de Datos y Garant{\'i}a de los Derechos Digitales (LOPDGDD) son
    aplicables porque los datos de contrataci{\'o}n incluyen informaci{\'o}n de car{\'a}cter
    personal (NIF de personas f{\'i}sicas que act{\'u}an como administradores o representantes).
    ProcureWatch opera exclusivamente sobre datos publicados en el marco de obligaciones
    legales de transparencia, lo que encuadra su tratamiento en la base jur{\'i}dica del
    art.~6.1.c RGPD (cumplimiento de una obligaci{\'o}n legal). El sistema no trata datos
    personales sensibles ni elabora perfiles individuales sin base legal espec{\'i}fica.

  \item[Reglamento (UE) 2024/1689 de Inteligencia Artificial (AI Act)] \cite{AIAct2024}
    Clasifica como \textbf{sistemas de alto riesgo} (Anexo III, secci{\'o}n~8) los sistemas
    que generen puntuaciones o clasificaciones de riesgo que puedan influir en decisiones
    sobre el acceso de empresas a contratos p{\'u}blicos. En consecuencia, ProcureWatch
    est{\'a} dise{\~n}ado con: (a)~trazabilidad completa del razonamiento del sistema para
    cada score generado; (b)~documentaci{\'o}n t{\'e}cnica conforme al art{\'i}culo~11; (c)~registro
    autom{\'a}tico de eventos (art{\'i}culo~12); y (d)~supervisi{\'o}n humana obligatoria: el
    sistema genera fichas de riesgo para revisi{\'o}n humana, sin capacidad de adoptar
    decisiones autom{\'a}ticas vinculantes (art{\'i}culo~14).
\end{description}

%=========================================================
\chapter{Conclusiones y Trabajo Futuro}
%=========================================================

[A desarrollar en la fase final del TFM, tras la implementaci{\'o}n y evaluaci{\'o}n del
prototipo sobre el dominio CPV~71.]

%=========================================================
% BIBLIOGRAF\'IA
%=========================================================
\bibliographystyle{apacite}
\bibliography{bibliografia}

\appendix

\chapter{Ap{\'e}ndice A: Subtipos CPV~71 incluidos en el an{\'a}lisis}

La Tabla~\ref{tab:cpv71} recoge los c{\'o}digos de la divisi{\'o}n CPV~71 incluidos en el
\textit{dataset} de validaci{\'o}n del prototipo, con su descripci{\'o}n oficial y la
presencia en el dataset BOE 2014--2024.

\begin{table}[ht]
\centering
\small
\caption{Subtipos CPV~71 relevantes para ProcureWatch Analytics}
\label{tab:cpv71}
\begin{tabular}{lll}
\toprule
\textbf{C{\'o}digo CPV} & \textbf{Descripci{\'o}n} & \textbf{Nivel} \\
\midrule
71000000 & Servicios de arquitectura, construcci{\'o}n, ingenier{\'i}a e inspecci{\'o}n & Divisi{\'o}n \\
71100000 & Servicios de arquitectura y conexos & Grupo \\
71200000 & Servicios de arquitectura y servicios conexos & Grupo \\
71300000 & Servicios de ingenier{\'i}a & Grupo \\
71311000 & Consultor{\'i}a en ingenier{\'i}a civil & Clase \\
71312000 & Consultor{\'i}a en ingenier{\'i}a de estructuras & Clase \\
71313000 & Consultor{\'i}a en ingenier{\'i}a medioambiental & Clase \\
71320000 & Servicios de dise{\~n}o t{\'e}cnico & Clase \\
71340000 & Servicios integrados de ingenier{\'i}a & Clase \\
71350000 & Servicios cient{\'i}ficos y t{\'e}cnicos en ingenier{\'i}a & Clase \\
71356000 & Servicios t{\'e}cnicos & Categor{\'i}a \\
71500000 & Servicios relacionados con la construcci{\'o}n & Grupo \\
71600000 & Servicios de pruebas, inspecci{\'o}n y control t{\'e}cnico & Grupo \\
71700000 & Servicios de supervisi{\'o}n e inspecci{\'o}n & Grupo \\
\bottomrule
\end{tabular}
\end{table}

\vspace{0.5em}
\noindent\textbf{Fuente:} \citeA{CPVCodes71}; elaboraci{\'o}n propia a partir del dataset BOE 2014--2024 \cite{MunozPla2026}.

\chapter{Ap{\'e}ndice B: Indicadores de red flags implementados (RF-01 a RF-15)}

La Tabla~\ref{tab:redflags} describe los 15 indicadores de red flags del motor
anal{\'i}tico de ProcureWatch, con el fundamento bibliogr{\'a}fico, el campo OCDS requerido
y el umbral de activaci{\'o}n propuesto.

\begin{longtable}{p{1.0cm}p{3.5cm}p{3.5cm}p{1.5cm}p{2.5cm}}
\caption{Indicadores de red flags RF-01 a RF-15}
\label{tab:redflags} \\
\toprule
\textbf{ID} & \textbf{Descripci{\'o}n} & \textbf{Campo OCDS} & \textbf{Umbral} & \textbf{Referencia} \\
\midrule
\endfirsthead
\multicolumn{5}{c}{\tablename\ \thetable{} (continuaci{\'o}n)} \\
\toprule
\textbf{ID} & \textbf{Descripci{\'o}n} & \textbf{Campo OCDS} & \textbf{Umbral} & \textbf{Referencia} \\
\midrule
\endhead
\bottomrule
\endfoot
RF-01 & Licitaci{\'o}n {\'u}nica & \texttt{bidsStatistics} & 1 oferta & \citeA{OCP2024RedFlags} \\
RF-02 & Concentraci{\'o}n de adjudicatario por organismo & \texttt{awards.suppliers} & $>$40\,\% & \citeA{Fazekas2020DIGIWHIST} \\
RF-03 & Desviaci{\'o}n importe estimado/adjudicado & \texttt{value} vs \texttt{award.value} & $>$20\,\% & \citeA{Coviello2023Italian} \\
RF-04 & Plazo de resoluci{\'o}n anormalmente corto & \texttt{tenderPeriod} & $<$10 d{\'i}as & \citeA{OCP2024RedFlags} \\
RF-05 & Recurrencia del mismo proveedor (negociado) & \texttt{procurementMethod} & $>$3 veces/a{\~n}o & \citeA{OIReScon2025} \\
RF-06 & Contrato menor cerca del umbral & \texttt{value.amount} & $>$90\,\% umbral & \citeA{OIReScon2025} \\
RF-07 & Publicaci{\'o}n en per{\'i}odo de baja visibilidad & \texttt{tender.tenderPeriod.startDate} & Festivo/dic. & \citeA{OCP2016RedFlags} \\
RF-08 & Modificaci{\'o}n contractual injustificada & \texttt{contracts.amendments} & $>$10\,\% & \citeA{OCP2024RedFlags} \\
RF-09 & Criterios subjetivos $>$49\,\% puntuaci{\'o}n & Texto pliego (NLP) & $>$49\,\% & \citeA{NLPLegalSurvey2024} \\
RF-10 & Concentraci{\'o}n en organismo-empresa & Grafo & $Q_{\text{local}}>0{,}30$ & \citeA{Traag2019Leiden} \\
RF-11 & Empresa con centralidad at{\'i}pica en red & Grafo betweenness & Top 1\,\% & \citeA{Wachs2025ML} \\
RF-12 & UTE con composici{\'o}n recurrente & Grafo + \texttt{suppliers} & $>$3 veces & Elaboraci{\'o}n propia \\
RF-13 & Plazo de ejecuci{\'o}n anormalmente corto & \texttt{contractPeriod} & $<$P5 dominio & \citeA{OCP2024RedFlags} \\
RF-14 & Ausencia de publicidad en TED (sobre umbral UE) & \texttt{tender.mainProcurementCategory} & Aplicable & \citeA{Directive2014} \\
RF-15 & Proveedor sin historial previo en dominio & \texttt{awards.suppliers.id} & Primera vez & \citeA{FalconCortes2021} \\
\end{longtable}

\end{document}
